%! Boxplots, Outliers, and Resistance — Unit 2, Lesson 2.3 — Warm-Up (Answer Key)
% ONE PAGE, blank and key. Three items at 12pt (\workrowsep 50pt, item spacing 22pt) so each
% item gets real handwriting room and still fits. Do not add a fourth item — the page is full.
% GRADUAL RELEASE: the warm-up activates Lesson 2.2's skills on the numbers the notes open with.
% Items 1-2 compute Pine Street Pizza's median, Q1, Q3 and IQR (the hook and rows 1, 2 and 4 of
% the notes reuse exactly these values: 25.5, 23, 28, 5). Item 3 turns counts into cumulative
% percents (25%, 50%, 75%) — row 3 of the notes graphs those same percents as the cumulative
% relative frequency graph.
% THE DATA: 25 30 22 55 27 20 24 28 23 26 -> in order 20 22 23 24 25 26 27 28 30 55;
%   median (25+26)/2 = 25.5; Q1 = med(20,22,23,24,25) = 23; Q3 = med(26,27,28,30,55) = 28;
%   IQR 5. Week: 10/40 = 25%, 20/40 = 50%, 30/40 = 75%.
\documentclass[12pt]{article}
\usepackage{saar-article}
\usepackage{saar-key}   % pulls in -boxes; do NOT also load -boxes
\newcommand{\TallMath}[1]{$\displaystyle #1\rule[-1.4em]{0pt}{3.2em}$}
% Word bank strip for every fill-in cluster (user request 2026-09-06). Defined per document;
% the key inherits it and prints the same strip, so heights match.
\newcommand{\wordbank}[1]{\par\vspace{2pt}\noindent\fcolorbox{goldacc}{goldbg}{\parbox{\dimexpr\linewidth-2\fboxsep-2\fboxrule\relax}{\small\textbf{Word bank:} #1}}\par\vspace{4pt}}
% Handwriting room inside every \begin{work} block. Moves the blank and the
% key together, so the two can never drift apart.
\setlength{\workrowsep}{50pt}

\begin{document}

\pageheader{Unit 2, Lesson 2.3}{Warm-Up --- Answer Key}
% NAMESTRIP: no \namedateperiod here — mirrors the blank. NO teachernote here —
% teacher prose lives in the lesson plan.

\vspace{0.15in}

\boxguard
\begin{notesbox}{Spiral Review --- Pine Street Pizza, Friday Night}
\normalsize
Pine Street Pizza timed its ten Friday-night deliveries. In the order they happened, in
minutes: \quad $25, \ 30, \ 22, \ 55, \ 27, \ 20, \ 24, \ 28, \ 23, \ 26$.

\begin{enumerate}[leftmargin=1.8em, itemsep=22pt, topsep=10pt]

  \item Write the ten times in order, then find the \textbf{median}.

        \par\vspace{6pt}
        In order: \ans{20, 22, 23, 24, 25, 26, 27, 28, 30, 55}

        \begin{work}
          \text{median} &= \dfrac{25 + 26}{2} \\
                        &= 25.5 \text{ minutes}
        \end{work}

  \item Find $Q_1$ and $Q_3$ --- the medians of the lower five and the upper five times ---
        then the \textbf{IQR}.

        \par\vspace{6pt}
        $Q_1 = $ \ans{23 minutes} \qquad $Q_3 = $ \ans{28 minutes}

        \begin{work}
          \text{IQR} &= 28 - 23 = 5 \text{ minutes}
        \end{work}

  \item Over the whole week the shop made $40$ deliveries. Write each count as a
        \textbf{percent} of all $40$.
        \wordbank{20\% \;\textbullet\; 25\% \;\textbullet\; 40\% \;\textbullet\; 50\% \;\textbullet\; 75\% \;\textbullet\; 80\%}

        \vspace{2pt}
        \renewcommand{\arraystretch}{2.1}
        \begin{tabularx}{\linewidth}{@{} X c p{3.4cm} @{}}
        \toprule
        \textbf{Deliveries that took\dots} & \textbf{Count} & \textbf{Percent of the 40} \\
        \midrule
        $20$ minutes or less & $10$ & \ans{$10/40 = 25\%$} \\
        $25$ minutes or less & $20$ & \ans{$20/40 = 50\%$} \\
        $30$ minutes or less & $30$ & \ans{$30/40 = 75\%$} \\
        \bottomrule
        \end{tabularx}

\end{enumerate}
\end{notesbox}

\end{document}
